\documentclass[11pt]{article}
\usepackage[a4paper,margin=27mm]{geometry}
\usepackage{amsmath,amssymb,amsthm,hyperref}
\newtheorem{theorem}{Theorem}
\newtheorem{lemma}[theorem]{Lemma}
\newtheorem{corollary}[theorem]{Corollary}
\newtheorem{remark}[theorem]{Remark}
\title{A quadratic-grid barrier for whole-copy insertion}
\author{Research notes, 30 September 2026}
\date{}
\begin{document}
\maketitle
\noindent\textbf{Status.} Complete proofs below. This is a barrier for a
specified construction framework, not a lower bound for arbitrary fair dice.
The quadratic scale of positive equidistant quadrature is classical;
the purpose here is a self-contained explicit bound and its consequence
for the dice insertion scheme.

\begin{theorem}[Positive quadrature requires a quadratic grid]
Let $r\ge2$ and suppose nonnegative weights on a subset of
$\{0,1/N,\ldots,1\}$ integrate every polynomial of degree at most $r$
exactly for the uniform probability measure on $[0,1]$. Put
$m=\lfloor(r-2)/2\rfloor$. Then
\[
 N\ge\frac{m^2+4m+6}{4}\ge\frac{r^2}{16}.
\]
The weights may be arbitrary real numbers; they need not be equal,
rational, or positive at every grid point.
\end{theorem}
\begin{proof}
Let $q(x)=P_m^{(0,3)}(2x-1)$, using the standard Jacobi normalization.
The following two norm identities are proved below:
\[
 A_m:=\int_0^1xq(x)^2\,dx=\frac{m^2+4m+6}{12},
 \qquad B_m:=\int_0^1x^2q(x)^2\,dx=\frac13.
\]
The degree of $f(x)=x(1/N-x)q(x)^2$ is $2m+2\le r$.
At every grid point, $f$ is nonpositive. Positivity and exactness
therefore imply
\[
 0\ge\sum_i w_i f(i/N)=\int_0^1f(x)\,dx=A_m/N-B_m.
\]
This yields the first bound. For even $r$, $m=r/2-1$; for odd $r$,
$m=(r-3)/2$. Direct substitution gives the second bound in both cases.
\end{proof}

Here is the promised verification of the Jacobi identities. Set
$J_k^{(b)}(x)=P_k^{(0,b)}(2x-1)$. Orthogonality gives
\[
 \int_0^1x^bJ_k^{(b)}(x)J_\ell^{(b)}(x)\,dx
   =\frac{\mathbf1_{k=\ell}}{2k+b+1}.
\]
The Jacobi contiguous relation, solved recursively, gives
\begin{align*}
 J_m^{(3)}
 &=\frac1{D_m}\sum_{k=0}^m(-1)^{m-k}
     (2k+3)(k+1)(k+2)J_k^{(2)},\\
 J_k^{(2)}
 &=\frac1{(k+1)(k+2)}\sum_{j=0}^k(-1)^{k-j}
     2(j+1)^2J_j^{(1)},
 \qquad D_m=(m+1)(m+2)(m+3).
\end{align*}
Thus
\[
 B_m=\frac1{D_m^2}\sum_{k=0}^m(2k+3)(k+1)^2(k+2)^2=\frac13;
\]
the sum telescopes because its summand is
\[
 \frac{(k+1)^2(k+2)^2(k+3)^2-k^2(k+1)^2(k+2)^2}{3}.
\]
Combining the two connection formulas yields
\[
 J_m^{(3)}=\frac2{D_m}\sum_{j=0}^m(-1)^{m-j}(j+1)^2
   \bigl((m+2)^2-(j+1)^2\bigr)J_j^{(1)}.
\]
Consequently
\[
 A_m=\frac2{D_m^2}\sum_{\ell=1}^{m+1}\ell^3
       \bigl((m+2)^2-\ell^2\bigr)^2
     =\frac{m^2+4m+6}{12},
\]
where the last equality follows by summing the powers $\ell^3,\ell^5,
\ell^7$ (or induction in $m$). The connection formulas follow directly
from
\[
 (2k+b+1)J_k^{(b)}=(k+b+1)J_k^{(b+1)}+kJ_{k-1}^{(b+1)}.
\]
Standard normalization and the contiguous relation are recorded at
\url{https://dlmf.nist.gov/18.3} and
\url{https://dlmf.nist.gov/18.9.E5}.

\begin{lemma}[Necessity of the insertion criterion]
Let $s$ be a permutation-fair string on $r$ letters, with arbitrary
multiplicities $c_1,\ldots,c_r$. Form $T$ from $N$ unchanged copies of
$s$, and insert runs of the new letter in gaps $i=0,\ldots,N$,
with nonnegative integer lengths $k_i$ and $K=\sum_i k_i>0$.
Then $T$ is permutation-fair if and only if $w_i=k_i/K$ gives a
degree-$r$ uniform quadrature on the grid $i/N$.
\end{lemma}
\begin{proof}
The partial-pattern formula for a fair string, and its powers, is
\[
 C_\pi(s^a)=\frac{a^{|\pi|}}{|\pi|!}
                  \prod_{\ell\in\pi}c_\ell.
\]
For a fixed full pattern of $T$ placing the new letter after $j$
old letters, its count therefore equals
\[
 \frac{KN^r\prod_{\ell=1}^r c_\ell}{r!}
  \sum_{i=0}^N w_i\binom rj(i/N)^j(1-i/N)^{r-j}.
\]
Its required fair count is
$KN^r\prod_\ell c_\ell/(r+1)!$. Thus fairness is equivalent to
the degree-$r$ Bernstein moment equations. Bernstein polynomials
span all polynomials of degree at most $r$, proving the claim.
\end{proof}

\begin{corollary}[An intrinsic $n\log n$ insertion cost]
Consider any sequence of constructions that starts from a one-letter
string and adds one new letter at a time by placing runs of it between
whole unchanged copies of the preceding string. Suppose every intermediate
string is permutation-fair. The oldest die in the final $n$-letter string
has at least
\[
 \prod_{r=2}^{n-1}\frac{r^2}{16}
 =\frac{((n-1)!)^2}{16^{n-2}}
 =\exp\bigl(2n\log n-O(n)\bigr)
\]
faces (the first displayed bound is merely weak when a factor is below $1$).
In particular, no choice of quadrature weights or improved denominator
estimate within this framework can give an $\exp(O(n))$ construction.
\end{corollary}
\begin{proof}
The $r$-to-$(r+1)$ insertion uses $N_r\ge r^2/16$ copies by the
preceding results. Every such insertion multiplies the oldest die's
face count by $N_r$. Multiply these inequalities and use Stirling's
formula.
\end{proof}

The same argument allows runs of zero length, hence omitted grid nodes,
and any grouping of the unchanged copies into larger blocks. It also
allows starting from an arbitrary fixed-size fair seed. It does not
apply when the old string is rearranged or locally changed at an insertion,
or when several letters are inserted jointly by a different mechanism.

\end{document}
